\section{Models (predictor types)}\label{sec:5}

\begin{disposition}
We will compare all predictor types we tried \emph{on the same folds}, from no learning at all to the final models (Table~\ref{tab:5.1}), and explain why some work better than others.
\end{disposition}

% Final text:


\begin{table}[H]
  \centering\small
  \caption{Predictor types on the same folds (A–D)}
  \label{tab:5.1}
  \adjustbox{max width=\linewidth}{\begin{tabular}{llllll}
\toprule
\textbf{} & \textbf{A (2020)} & \textbf{B (2021)} & \textbf{C (2022)} & \textbf{D (2019)} & \textbf{mean} \\
\midrule
Rule: price − upstream water value (+7 d) & 0.8559 & 0.8488 & 0.9115 & 0.8629 & 0.8698 \\
Relaxed LP alone: unit on & 0.9318 & 0.9483 & 0.9554 & 0.9291 & 0.9411 \\
Relaxed LP alone: marginal profit & 0.951 & 0.961 & 0.9595 & 0.9568 & 0.9571 \\
LightGBM without LP features (exp 009) & 0.9667 & 0.9837 & 0.9834 & 0.9817 & 0.9789 \\
XGBoost without LP features (exp 012) & 0.965 & 0.9833 & 0.983 & 0.9806 & 0.978 \\
LightGBM, one model per unit (exp 028f) & 0.9808 & 0.9917 & 0.9872 & 0.9804 & 0.985 \\
LightGBM, pooled (exp 029, final 1) & 0.9868 & 0.9941 & 0.9922 & 0.9919 & 0.9912 \\
XGBoost, pooled (exp 030, final 2) & 0.986 & 0.994 & 0.9922 & 0.9916 & 0.991 \\
Blend 50/50 of exp 029 and 030 & 0.9867 & 0.9942 & 0.9924 & 0.9919 & 0.9913 \\
\bottomrule
\end{tabular}
}
\end{table}

\begin{disposition}
\emph{Note:} exp 028f (one model per unit) uses the LP features of exp 024, the previous LP version; compare it with exp 024 (0.9865), not with exp 029.
\end{disposition}

% Final text:


\subsection{Rules without learning (baselines)}\label{sec:5.1}

\begin{disposition}
\begin{enumerate}
\def\labelenumi{(\alph{enumi})}
\tightlist
\item
  The single economic rule ``price − upstream water value at +7 d'' ranks all unit-hours by one number. (b) The relaxed LP alone, ranking by its fractional ``unit on'' or by its marginal profit. Including these shows how much comes from the physics and how much from the machine learning: the LP alone is already strong, and the GBDT on top closes most of the remaining gap.
\end{enumerate}
\end{disposition}

% Final text:


\subsection{Gradient-boosted trees: LightGBM (final model 1)}\label{sec:5.2}

\begin{disposition}
Why this model type: the data is tabular (one row per case × unit × hour, about 6.9 million rows); trees capture interactions such as ``the partner plant runs \emph{and} the margin is positive'' without us specifying them; they are robust to features on different scales and to redundant features; they handle missing values (the partner features are empty for most units); and they train in under a minute. With about 2900 cases, tree ensembles are the usual strong choice for tabular data (Grinsztajn et al., 2022). We use one \emph{pooled} model for all 14 units with the unit ID as a categorical feature.
\end{disposition}

% Final text:


\subsection{XGBoost (final model 2)}\label{sec:5.3}

\begin{disposition}
Same features, settings matched to LightGBM (leaf-wise growth, same number of leaves and trees). It scores the same as LightGBM (0.9910 vs 0.9912), and its test predictions correlate 0.997 with LightGBM's. Explanation: both are gradient-boosted trees, and with these features the information, not the learning algorithm, sets the limit. We still chose it as the second submission because it is a different implementation, which gives some insurance on the private leaderboard, and its public score was slightly higher.
\end{disposition}

% Final text:


\subsection{One model per unit versus one pooled model (exps 010, 028f)}\label{sec:5.4}

\begin{disposition}
Separate models for each unit were worse (−0.0016 and −0.0014). Explanations: (1) the economics are shared by all units, so a pooled model learns them from 14× more data; (2) separate models produce probabilities on separate scales, while micro-AUC ranks all units in one list, so a common scale helps directly.
\end{disposition}

% Final text:


\subsection{Blends (exp 020 and LightGBM + XGBoost on the final features)}\label{sec:5.5}

\begin{disposition}
Averaging LightGBM and XGBoost gave +0.0000 to +0.0001. Explanation: blending only helps when models make \emph{different} errors, and these two make nearly the same errors because they see the same features.
\end{disposition}

% Final text:


\subsection{Integer optimisation (MILP) instead of the LP}\label{sec:5.6}

\begin{disposition}
Solving the LP with true on/off decisions is closer to SHOP but gave \emph{worse} features (units-on AUC 0.911 vs 0.945). Explanation: AUC rewards knowing \emph{how close} a decision is; the fractional LP value carries that information, a hard 0/1 decision throws it away.
\end{disposition}

% Final text:


\subsection{Post-processing we tried}\label{sec:5.7}

\begin{disposition}
Smoothing predictions over neighbouring hours (worse at every window size, because the labels switch sharply and the LP timing is precise); enforcing the plant priority order (no effect, the model already respects it); averaging predictions over overlapping cases (worse, because neighbouring cases are solved independently and disagree on 12\% of hours).
\end{disposition}

% Final text:


\subsection{Neural networks (not trained)}\label{sec:5.8}

\begin{disposition}
We will be honest that we did not train a neural network, and explain why: with about 2900 weeks of data, tree models on tabular features are the safer choice; SINTEF's paper shows that neural networks on this dataset mainly add a small correction on top of the relaxed LP's solution; and our time went into the LP, which was where the gains were. A 1D CNN over the LP's hourly outputs would be the natural next model to try.
\end{disposition}

% Final text:

